\documentclass[11pt,a4paper]{article}
\usepackage[T1]{fontenc}
\usepackage[utf8]{inputenc}
\usepackage{lmodern}
\usepackage[a4paper,margin=25mm,headheight=15pt]{geometry}
\usepackage{amsmath,amssymb,amsthm,mathtools}
\usepackage{microtype,booktabs,longtable,array,tabularx}
\usepackage[dvipsnames]{xcolor}
\usepackage{listings,fancyhdr}
\usepackage[hyphens]{url}
\usepackage{hyperref}
\hypersetup{colorlinks=true,linkcolor=MidnightBlue,citecolor=MidnightBlue,
  urlcolor=MidnightBlue,pdftitle={Exact Functional Decomposition of Polynomials with Algebraic Coefficients},
  pdfauthor={Prepared for Vladimir Reshetnikov},pdfsubject={Theory, exact certificates, and Wolfram Language and Python implementations}}
\pagestyle{fancy}
\fancyhf{}
\fancyhead[L]{\small Exact functional decomposition}
\fancyhead[R]{\small Algebraic coefficients}
\fancyfoot[C]{\thepage}
\renewcommand{\headrulewidth}{0.3pt}
\setlength{\parskip}{4pt}
\setlength{\parindent}{1.2em}
\setlength{\emergencystretch}{2em}
\numberwithin{equation}{section}
\newtheorem{theorem}{Theorem}[section]
\newtheorem{lemma}[theorem]{Lemma}
\newtheorem{proposition}[theorem]{Proposition}
\newtheorem{corollary}[theorem]{Corollary}
\theoremstyle{definition}
\newtheorem{definition}[theorem]{Definition}
\newtheorem{example}[theorem]{Example}
\theoremstyle{remark}
\newtheorem{remark}[theorem]{Remark}
\newcommand{\Q}{\mathbb{Q}}
\newcommand{\Qbar}{\overline{\mathbb{Q}}}
\newcommand{\lc}{\operatorname{lc}}
\newcommand{\ord}{\operatorname{ord}}
\newcommand{\coeff}[2]{[#1]#2}
\newcommand{\code}[1]{\texttt{#1}}
\lstdefinestyle{wl}{basicstyle=\ttfamily\footnotesize,
  columns=fullflexible,keepspaces=true,breaklines=true,breakatwhitespace=false,
  showstringspaces=false,frame=single,framerule=0.25pt,
  rulecolor=\color{black!25},xleftmargin=5pt,xrightmargin=5pt,
  aboveskip=8pt,belowskip=8pt,tabsize=2,
  morecomment=[s]{(*}{*)},commentstyle=\color{black!60},
  morestring=[b]",stringstyle=\color{MidnightBlue}}
\lstset{style=wl}
\begin{document}

\title{Exact functional decomposition of polynomials\\
with algebraic coefficients\\[4pt]
\large A unified theory, algorithms, certificates, and implementations}
\author{Prepared for Vladimir Reshetnikov\\
\small in response to Mathematica StackExchange question 206618}
\date{September 2026}
\maketitle
\begin{abstract}
A polynomial whose coefficients contain radicals or algebraic \code{Root} objects can be
functionally decomposed without factoring its derivative.  In characteristic zero, normalizing
an inner component to be monic with constant term zero removes the affine ambiguity: its degree
alone determines a unique candidate from the leading coefficients of the input.  A finite
formal-root recurrence constructs it in the coefficient field, and exact polynomial division
supplies an acceptance or rejection certificate.

This article synthesizes three technical reports and the accompanying Wolfram Language and
Python implementations.  We give self-contained proofs of the fixed-degree algorithm, descent
through extensions, complete normalized-chain enumeration, and the divisibility poset of all
right components.  We distinguish two outer-recovery methods and their negative certificates,
derive explicit low-degree criteria, and work through radical, nested, quintic \code{Root},
complex, monomial, and Chebyshev examples.  Arithmetic complexity, algebraic representation,
the derivative-based answer in the original discussion, the differences among the reports,
and current native and independent validation are treated separately.
\end{abstract}
{\setlength{\parskip}{0pt}\small\tableofcontents}
\clearpage

\section{Problem, scope, and main conclusion}
\label{sec:intro}
The motivating question asks whether \code{Decompose} can be made to work on
polynomials with radical coefficients, and presents a derivative-factorization
implementation as an alternative \cite{SE}. The central distinction is between
\emph{functional decomposition}
\begin{equation}
 p(x)=f(h(x)),\qquad \deg f>1,\quad \deg h>1,
 \label{eq:problem}
\end{equation}
and ordinary multiplicative factorization. They are different problems.
Irreducibility under multiplication neither implies nor is implied by
indecomposability under composition.

For instance, $x^4+1=(t^2+1)\circ x^2$ is functionally decomposable even though
it is irreducible over $\Q$. Conversely, $x^4+x$ is multiplicatively reducible
but functionally indecomposable in characteristic zero, as the algorithm below
immediately proves. An \code{Extension} option appropriate for multiplicative
factorization is therefore not the fundamental missing ingredient.

Let
\begin{equation}
 p(x)=\sum_{i=0}^{n}a_i x^i,\qquad a_n\ne0,
 \qquad K=\Q(a_0,\ldots,a_n)\subseteq\Qbar.
 \label{eq:input}
\end{equation}
Because there are finitely many algebraic coefficients, $K/\Q$ is a finite
extension. The input expressions choose particular algebraic numbers, not
unspecified conjugacy classes. A radical expression, an algebraic
\code{Root}, and an \code{AlgebraicNumber} can denote the same element of $K$.
Their outward forms should not determine the answer.

\paragraph{The central result.}
For every proper divisor $d$ of $n$, there is at most one pair $(f,h)$ satisfying
\eqref{eq:problem} with $\deg h=d$, $h$ monic, and $h(0)=0$. Its only possible
$h$ is computed from the leading $d$ coefficients of $p$ by rational
operations. A subsequent exact subalgebra-membership test either constructs
$f$ or certifies its nonexistence. No polynomial factorization, solving of
coefficient equations, extraction of new algebraic numbers, or search through
field extensions is needed.

This is a characteristic-zero realization of the classical factorization-free
approach to polynomial decomposition, particularly the work of Kozen and
Landau \cite{KL}. Faster and more general algorithmic treatments exist,
including the tame-case work of von zur Gathen \cite{Gathen}. The elementary
algorithm here is chosen for transparent proofs, straightforward exact
implementation, and verifiable failure certificates, not as a claim to a new
asymptotically optimal decomposition algorithm.

\paragraph{What is and is not claimed about the built-in function.}
The source discussion is from 2019. The currently retrieved Wolfram
documentation illustrates symbolic and complex coefficients; it does not
state an integer-only mathematical domain for \code{Decompose} \cite{Decompose}.
The historical behavior reported in the question must not be promoted to a
verified claim about every current Wolfram version. The native validation
reported below concerns the delivered package. It is an independent implementation;
it does not modify \code{System`Decompose} or depend on its behavior.

\paragraph{Scope of the delivered code.}
The theory applies to any characteristic-zero field with exact arithmetic.
The software intentionally accepts the narrower, decidable input class of
explicit exact algebraic coefficients recognizable by their respective
exact algebraic-number backends. It does not accept approximate coefficients, symbolic
parameters, transcendental coefficients, rational functions in $x$, or
finite-characteristic arithmetic. These exclusions are input contracts,
not assertions that broader mathematical variants are impossible.

\section{Normalization and the coefficient field}
\label{sec:normalization}
\begin{definition}
A nonconstant polynomial $h\in K[x]$ is \emph{normalized} if it is monic and
$h(0)=0$. A decomposition is \emph{nontrivial} if both component degrees exceed
one. A \emph{complete chain} is a composition of polynomials of degree at
least two, each of which is functionally indecomposable.
\end{definition}

Suppose $p=f\circ g$ and write $b=\lc(g)\ne0$, $c=g(0)$. Set
\begin{equation}
 h(x)=\frac{g(x)-c}{b},\qquad F(t)=f(bt+c).
 \label{eq:normalize}
\end{equation}
Then $p=F\circ h$, and $h$ is normalized. This replaces the arbitrary affine
freedom at a composition interface by a canonical representative.

More explicitly, if $\ell(t)=bt+c$, the pairs
$(f,g)$ and $(f\circ\ell^{-1},\ell\circ g)$ describe the same composition.
Without removing this freedom there are generally infinitely many pairs even
for a fixed degree pattern. We enumerate representatives modulo precisely
these affine interface changes; we do not enumerate arbitrary choices of
linear factors.

In a normalized complete chain
\begin{equation}
 p=f_1\circ f_2\circ\cdots\circ f_r,
 \label{eq:chain}
\end{equation}
every component except $f_1$ is monic with constant term zero. There is no
normalization requirement on $f_1$: it must carry the input's leading
coefficient and constant term. This convention also explains why a correct
answer may look different from the radical expression in the question.

\begin{remark}
An algebraic number's selected embedding matters. For example, replacing
$\sqrt2$ by $-\sqrt2$ changes the input polynomial in general. What is
representation-independent is replacing an expression by another exact
expression for the \emph{same} algebraic number. No \code{PowerExpand}-style
branch assumptions are legitimate here.
\end{remark}

\section{The unique right-component candidate}
\label{sec:candidate}
Fix $1<d<n$ with $d\mid n$ and set $m=n/d$. Seek a normalized component
\begin{equation}
 h(x)=x^d+b_1x^{d-1}+\cdots+b_{d-1}x.
 \label{eq:h}
\end{equation}
If $p=f(h)$, monicity of $h$ forces $\lc(f)=a_n$. All terms of
$f(h)-a_nh^m$ have degree at most $(m-1)d=n-d$. Consequently,
\begin{equation}
 [x^{n-k}]h(x)^m=\frac{a_{n-k}}{a_n},
 \qquad 1\le k\le d-1.
 \label{eq:topcongruence}
\end{equation}
These are only $d-1$ equations, exactly as many as the unknown coefficients
of the normalized $h$.

\begin{lemma}[Triangular recovery]
\label{lem:triangular}
Over a characteristic-zero field, equations \eqref{eq:topcongruence} have a
unique solution $b_1,\ldots,b_{d-1}\in K$. They can be solved successively by
\begin{equation}
 b_k=\frac1m\left(\frac{a_{n-k}}{a_n}-
 [x^{n-k}]\left(x^d+\sum_{j=1}^{k-1}b_jx^{d-j}\right)^m\right).
 \label{eq:triangular}
\end{equation}
\end{lemma}
\begin{proof}
View a choice of $b_jx^{d-j}$ from one of the $m$ factors as a degree deficit
of $j$. A contribution of total deficit $k$ either chooses $b_k$ once and
leading terms in every other factor, or uses only deficits smaller than $k$.
The coefficient in the first case is $mb_k$. No $b_j$ with $j>k$ can contribute.
Thus the coefficient equation is linear in the new $b_k$, with nonzero
coefficient $m$. Induction proves existence and uniqueness. Every step uses
arithmetic in $K$ alone.
\end{proof}

\begin{corollary}[Uniqueness for a fixed degree]
\label{cor:unique}
For fixed $d$, at most one normalized right component can occur. Once it
occurs, its outer component is also unique.
\end{corollary}
\begin{proof}
The first assertion follows from Lemma~\ref{lem:triangular}. Substitution
$K[t]\to K[x]$, $f\mapsto f(h)$, is injective for nonconstant $h$, since a
nonzero polynomial of degree $r$ maps to one of degree $rd$.
\end{proof}

\subsection{A quadratic-cost formal-root recurrence}
Repeatedly expanding powers in \eqref{eq:triangular} is conceptually simple
but unnecessary. Put
\begin{equation}
 S(t)=1+\sum_{i=1}^{d-1}s_i t^i\pmod{t^d},
 \qquad s_i=\frac{a_{n-i}}{a_n}.
 \label{eq:S}
\end{equation}
There is a unique $U(t)=1+\sum_{k=1}^{d-1}u_k t^k\pmod{t^d}$ with
$U(t)^m=S(t)\pmod{t^d}$. It is the finite formal series with constant term
one, customarily written $S(t)^{1/m}$. It is not a numerical choice of a
complex $m$th-root branch. Its coefficients are in $K$.

Differentiating $U^m=S$ in formal power series gives
$mSU'=US'$. Comparing coefficients of $t^{k-1}$ yields
\begin{equation}
 \boxed{\quad
 u_0=1,\qquad
 u_k=\frac{1}{mk}\sum_{i=1}^{k}
       \bigl((m+1)i-mk\bigr)s_i u_{k-i},
 \quad 1\le k<d.\quad}
 \label{eq:recurrence}
\end{equation}
For completeness, the left side before isolating $u_k$ is
$mku_k+m\sum_{i=1}^{k-1}(k-i)s_i u_{k-i}$; the right side is
$\sum_{i=1}^k i s_i u_{k-i}$. Their difference is precisely
\eqref{eq:recurrence}. The candidate is
\begin{equation}
 h_d(x)=x^d+u_1x^{d-1}+\cdots+u_{d-1}x.
 \label{eq:hd}
\end{equation}
The first three coefficients illustrate the rational nature of the construction:
\begin{align}
 u_1&=\frac{s_1}{m},\notag\\
 u_2&=\frac{s_2}{m}-\frac{m-1}{2m^2}s_1^2,\notag\\
 u_3&=\frac{s_3}{m}-\frac{m-1}{m^2}s_1s_2
       +\frac{(m-1)(2m-1)}{6m^3}s_1^3.
 \label{eq:firstcoeffs}
\end{align}
Only terms with indices $k<d$ are needed. The inner constant term is fixed
at zero rather than inferred from a formal-root coefficient of order $d$.

\subsection{The Laurent-series interpretation}
At infinity, write
\[
 \left(p(x)/a_n\right)^{1/m}
 =x^d\left(1+\frac{a_{n-1}}{a_n}x^{-1}+\cdots\right)^{1/m},
\]
using the branch whose leading coefficient is one. The positive-degree
polynomial part is exactly $h_d$. If $p=f(h)$ with
$f(t)=a_nt^m+c_{m-1}t^{m-1}+\cdots$, then
\[
 \left(p/a_n\right)^{1/m}
 =h+\frac{c_{m-1}}{ma_n}+O(h^{-1}).
\]
Thus the full polynomial part may contain a nonzero constant, but subtracting
that constant gives the normalized $h$. This provides an approximate-root
interpretation of the same exact finite computation. It does \emph{not}
justify accepting the candidate without testing the lower coefficients of
$p$.

\subsection{Why characteristic zero is explicitly required}
The triangular construction divides by $m$; the efficient recurrence also
divides by $k$. Characteristic zero makes both safe. In characteristic two,
\[
 x^4+x^2=(t^2+t)\circ x^2=t^2\circ(x^2+x),
\]
and the two right components are distinct normalized polynomials of the same
degree. Therefore the fixed-degree uniqueness theorem itself can fail in wild
characteristic. Even when $m$ is invertible in positive characteristic, the
particular recurrence \eqref{eq:recurrence} cannot simply be used if some
$k<d$ is not invertible. The delivered code makes no finite-field claim.

\section{Exact membership in the subalgebra \texorpdfstring{$K[h]$}{K[h]}}
\label{sec:digits}
The candidate equations concern only the top of $p$. To test the whole
polynomial, regard $h$ as a polynomial ``base.'' No unknown outer coefficients
or general-purpose solver are required.

\begin{theorem}[Polynomial-base expansion]
\label{thm:digits}
Let $h\in K[x]$ be monic of degree $d\ge1$. Every $p\in K[x]$ has a unique
finite expansion
\begin{equation}
 p(x)=\sum_{j\ge0}r_j(x)h(x)^j,\qquad \deg r_j<d.
 \label{eq:digits}
\end{equation}
Moreover, $p\in K[h]$ if and only if every $r_j$ is constant.
\end{theorem}
\begin{proof}
Repeated Euclidean division by the monic $h$ gives
$q_j=hq_{j+1}+r_j$, starting at $q_0=p$. Quotient degree drops by $d$ whenever
it is at least $d$, so the process terminates. Substitution gives
\eqref{eq:digits}.

For uniqueness, suppose a nonzero finite combination
$\sum r_jh^j$ vanishes, and let $J$ be the largest index with $r_J\ne0$.
Its leading degree lies in the interval $[Jd,(J+1)d-1]$, whereas every smaller
index contributes a degree below $Jd$. Cancellation of its leading term is
impossible. Hence all digits in a zero combination are zero.

If the digits are constants, the expansion itself defines an outer polynomial
in $K[t]$. Conversely, an expression $p=\sum c_jh^j$ with $c_j\in K$ is a
valid digit expansion. Uniqueness forces $r_j=c_j$.
\end{proof}

Equivalently, $K[x]$ is a free $K[h]$-module with basis
$1,x,\ldots,x^{d-1}$. The nonconstant digit coefficients are the coordinates
that prevent membership in its rank-one submodule $K[h]\cdot1$.

For $\deg p=md$, the digit list is $r_0,\ldots,r_m$ and $r_m=a_n$.
The outer candidate used by the implementation is
\begin{equation}
 F_d(t)=\sum_{j=0}^{m}r_j(0)t^j.
 \label{eq:outercandidate}
\end{equation}
It is an actual outer component precisely when all the digits are constant.
On rejection the package reports both the nonzero residual $p-F_d(h_d)$ and
one nonconstant digit coefficient. The latter identifies the exact failure,
not merely an unsuccessful solver run.

\begin{theorem}[Necessary and sufficient fixed-degree algorithm]
\label{thm:fixed}
Let $K$ have characteristic zero and let $d$ be a proper divisor of
$n=\deg p$. Construct $h_d$ by \eqref{eq:recurrence}--\eqref{eq:hd}, and
compute the digits of $p$ in base $h_d$. There exists a decomposition with
right degree $d$ if and only if every digit is constant. On success,
$(F_d,h_d)$ is the unique normalized pair.
\end{theorem}
\begin{proof}
Any normalized right component of degree $d$ must equal $h_d$ by
Lemma~\ref{lem:triangular}. Its existence is therefore equivalent to
$p\in K[h_d]$, which Theorem~\ref{thm:digits} decides exactly. The outer
component on success is \eqref{eq:outercandidate}; its degree is $m$ because
$r_m=a_n\ne0$. Corollary~\ref{cor:unique} supplies uniqueness.
\end{proof}

\paragraph{A small negative example.}
For $p=x^6+x$, both proper right degrees are $2$ and $3$. In either case the
leading coefficients force $h_d=x^d$. The digit $r_0=x$ is nonconstant, so
both degrees are impossible. This proves functional indecomposability even
though six is composite. The proof works over every characteristic-zero
extension of the coefficient field.

\subsection{A second exact recovery method: descending pivots}

Report~\cite{r2} instead recovers an outer polynomial by descending coefficient
cancellation.  Start with \(R=p\), and for \(j=m,m-1,\ldots,0\) set
\[
 c_j=[x^{jd}]R,\qquad R\leftarrow R-c_jh^j.
\]
The final representation is
\begin{equation}\label{eq:pivot-certificate}
 p=\widetilde F(h)+\widetilde R,\qquad
 \widetilde F=\sum_{j=0}^m c_jt^j,\qquad
 [x^{jd}]\widetilde R=0\quad(0\le j\le m).
\end{equation}
It is unique: the leading degree of a nonzero difference of outer polynomials
after substitution is divisible by \(d\), whereas a difference of two such
residuals has zero coefficients at every divisible degree.  Thus
\(p\in K[h]\) if and only if \(\widetilde R=0\).

The pivot and digit methods agree on acceptance and on every successful
outer component.  Their outer candidates and residuals can differ on rejection.
For example, take
\[
 h=x^2+x,\qquad
 p=h^3+xh=x^6+3x^5+3x^4+2x^3+x^2.
\]
This \(h\) is the forced degree-two candidate.  The base-\(h\) digits
are \((0,x,0,1)\), so the digit-constant construction gives
\(F=t^3\), \(R=x^3+x^2\).  Descending pivots instead give
\(\widetilde F=t^3+t\), \(\widetilde R=x^3-x\).
Both are correct negative tests, but they are different certificate schemas.
The live interface uses the digit convention of reports~\cite{r1,r3}.

In particular, the nonzero identity
\(p=G(h)+R\) alone is insufficient as a rejection certificate.
A pivot checker must verify the zero-pivot conditions
\eqref{eq:pivot-certificate}, while either checker must also verify that
\(h\) is the uniquely forced normalized candidate.

\section{Descent, functoriality, and decomposition loci}
\label{sec:descent}
\begin{theorem}[Descent through field extensions]
\label{thm:descent}
Let $L/K$ be any extension of characteristic-zero fields and $p\in K[x]$.
For any $m,d>1$, a decomposition of degree pattern $(m,d)$ exists over $L$
if and only if it exists over $K$. Every normalized pair over $L$ already
has both components in $K[x]$.
\end{theorem}
\begin{proof}
The forward implication is the only one to prove. Normalize a pair over $L$
by \eqref{eq:normalize}. Its right component is the unique solution of the
triangular coefficient system. Formula \eqref{eq:triangular} shows inductively
that all its coefficients lie in $K$. Monic division of $p\in K[x]$ by this
$h\in K[x]$ then produces digits in $K[x]$. Since they are constant over $L$,
they are elements of $K$, and these are the coefficients of the outer
component. The argument neither assumes that $L/K$ is algebraic nor uses a
Galois closure.
\end{proof}

The descent mechanism is a standard consequence of the leading-coefficient
method; related formulations appear explicitly in Wyman and Zieve
\cite{WZ}. Its relevance here is particularly strong:
\begin{corollary}
For the algebraic input field $K$ in \eqref{eq:input}, decomposition over $K$
and over $\Qbar$ gives the same normalized pairs and chains. Enlarging
\code{Extension} cannot create a missing normalized decomposition.
\end{corollary}

This conclusion is special to functional decomposition in characteristic
zero. It contrasts with ordinary factorization, where $x^2-2$ is irreducible
over $\Q$ but splits after adjoining $\sqrt2$. It also removes any need to
decide whether a coefficient's minimal polynomial is solvable by radicals.
A quintic \code{Root} is an entirely ordinary exact coefficient for this
algorithm.

\begin{proposition}[Compatibility with embeddings]
If $\sigma:K\hookrightarrow L$ is a field embedding, then candidate
construction, digit expansion, and normalized decomposition commute with
$\sigma$. In particular, conjugate input polynomials have the same accepted
right degrees and the same complete degree patterns.
\end{proposition}
\begin{proof}
Every construction uses finite additions, multiplications, and divisions by
nonzero leading coefficients or nonzero integers. These operations commute
with a field embedding. Nonzero coefficients stay nonzero, so the digit
obstructions and acceptance decisions agree.
\end{proof}

\subsection{A coefficient-space interpretation}
The finite tests also describe exactly where decompositions occur in families.
Fix $n=md$ and regard $a_0,\ldots,a_n$ as coordinates, with $a_n$ invertible.
The coefficients of $h_d$ and all its base digits are regular functions in
\[
 K[a_0,\ldots,a_n,a_n^{-1}]
\]
when the base field has characteristic zero. In other words, their only
variable denominators are powers of the leading coefficient; integer
constants in denominators are harmless.

\begin{proposition}[Fixed-degree decomposition locus]
\label{prop:locus}
Within the space of degree-$n$ polynomials, the locus admitting right degree
$d$ is the common zero set of the nonconstant digit coefficients. It is
isomorphic to the parameter space of arbitrary degree-$m$ outer polynomials
and normalized degree-$d$ inner polynomials. In particular it is smooth and irreducible. Its dimension is $m+d$ and its
codimension in the $(n+1)$-dimensional input coefficient space is
$(m-1)(d-1)$.
\end{proposition}
\begin{proof}
Theorem~\ref{thm:fixed} gives precisely the stated equations. The composition
map from parameters $(f,h)$ is polynomial. It is injective by normalized
uniqueness, and its inverse on the locus is the regular candidate-and-digit
construction just described. There are $m+1$ outer coefficients, with the
leading one nonzero, and $d-1$ free inner coefficients. This parameter space is the product of a punctured affine line and an affine space, so it is smooth and irreducible. Hence the dimension
is $(m+1)+(d-1)=m+d$. Subtracting from $n+1=md+1$ gives
$md+1-m-d=(m-1)(d-1)$.
\end{proof}

This description is not an extra symbolic-parameter feature in the package;
it is a mathematical explanation of specialization. As long as degree is
preserved, the formulas specialize correctly, but exceptional parameter
values may cause all obstructions to vanish. For example, $x^6+\lambda x$
is indecomposable for nonzero $\lambda$ and decomposes when $\lambda=0$.
Therefore no numerical ``generic behavior'' test can replace the exact
coefficient tests.

\section{One complete chain and all complete chains}
\label{sec:chains}
Let
\[
 D(p)=\{d:1<d<\deg p,\ d\mid\deg p,\
                   p\in K[h_d]\}.
\]
Testing all proper divisors computes every normalized nontrivial pair. There
are at most $\tau(n)-2$ such pairs for degree $n\ge2$, where $\tau(n)$ is the
number of positive divisors of $n$. Testing only prime right degrees is not
complete: $(x^4+x)^2$ has the atomic degree-four right component $x^4+x$ and
has no right component of degree two. Likewise, there is no justification for
restricting the search to $d\le\sqrt n$.

\begin{lemma}[Minimal right degree gives an atom]
\label{lem:minimal}
If $D(p)$ is nonempty and $d$ is its smallest member, then $h_d$ is
indecomposable.
\end{lemma}
\begin{proof}
If $h_d=u\circ v$ with both degrees greater than one, then
$p=(F_d\circ u)\circ v$. Normalizing $v$ yields a right component of $p$ of
degree $1<\deg v<d$, contradicting minimality.
\end{proof}

\subsection{Deterministic single-chain algorithm}
Test the proper divisors in increasing order. At the first success, record
$h_d$ as the rightmost component and repeat on $F_d$. Stop when no proper
degree succeeds. The outer degree decreases by at least a factor of two at
each step. The resulting list is returned outermost first.

\begin{theorem}[Completeness of the single-chain algorithm]
The algorithm terminates and returns a complete normalized chain for every
input of degree at least two. If it returns a singleton, the input is
indecomposable over every characteristic-zero extension of $K$.
\end{theorem}
\begin{proof}
Termination follows from strictly decreasing positive degree. Each recorded
right component is indecomposable by Lemma~\ref{lem:minimal}. The final outer
component is indecomposable because every possible right degree has been
rejected by Theorem~\ref{thm:fixed}. All recorded identities are exact, so
substituting them reconstructs $p$. Descent gives the extension-field claim.
\end{proof}

\subsection{Enumeration without repeated reassociations}
For all complete chains, it is tempting to split recursively on every pair
and on both sides. That repeats the same chain through different binary
parenthesizations. The package instead selects the last atom first:
\begin{equation}
 \mathcal C(p)=
 \begin{cases}
  \{(p)\}, & D(p)=\varnothing,\\[2pt]
  \displaystyle\bigcup_{\substack{d\in D(p)\\h_d\text{ indecomposable}}}
       \{(c_1,\ldots,c_r,h_d):
                          (c_1,\ldots,c_r)\in\mathcal C(F_d)\},
       & D(p)\ne\varnothing.
 \end{cases}
 \label{eq:enumeration}
\end{equation}
The implementations memoize pair searches and atomicity decisions locally
within one call. All-chain traversal is incremental; a requested output limit
can stop enumeration without first constructing all remaining chains.

\begin{theorem}[Complete enumeration modulo affine interfaces]
\label{thm:enumeration}
Formula \eqref{eq:enumeration} returns exactly all normalized complete chains,
with no duplicates caused by binary reassociation. In particular, a specified
ordered list of component degrees admits at most one normalized chain.
\end{theorem}
\begin{proof}
In any normalized complete chain, the final component is indecomposable and
has some degree $d\in D(p)$. By fixed-degree uniqueness, it must equal $h_d$,
and the rest of the chain composes to the unique $F_d$. Induct on degree for
the outer chain. This proves both inclusion directions and uniqueness for a
fixed ordered degree list. Different selected final degrees cannot produce
the same chain, and within one final degree the induction hypothesis rules
out duplicates. Strict degree reduction proves finiteness.
\end{proof}

A complete chain has at most $\lfloor\log_2 n\rfloor$ components, but the
number of different chains need not be small. Enumerating all of them is
necessarily output-sensitive. The interface intentionally separates one
chain, all pairs, and all chains so that a user need not request the largest
output to solve the ordinary decomposition problem.

\paragraph{Constants and affine polynomials.}
For software convenience, degree-zero and degree-one inputs return the
singleton list containing the input, and the all-chains function returns a
list containing that singleton. These are conventions, not complete chains
of nonlinear atoms in the definition above. The certificate metadata marks
indecomposability as not applicable below degree two. The zero polynomial's
reported degree is $-\infty$.

\section{The finite right-component poset}\label{sec:poset}
Assume throughout this section that $n=\deg p\ge2$.
The degree uniqueness theorem has a stronger consequence that is useful for understanding the whole decomposition structure. Include the two trivial normalized components
\[
 h_1=x,\qquad h_n=\frac{p-p(0)}{a_n},
\]
and let $S(p)$ be their degrees together with all successful proper right degrees.

\begin{theorem}[Divisibility is the right-component order]\label{thm:poset}
For $d,e\in S(p)$, there exists a polynomial $q$ satisfying $h_e=q\circ h_d$ if and only if $d\mid e$. If it exists, $q$ is unique, monic, has zero constant term, and has degree $e/d$.
\end{theorem}
\begin{proof}
The forward implication follows from multiplication of degrees, and uniqueness follows from injectivity of substitution into a nonconstant polynomial. Leading coefficients and constant terms give the normalization of $q$.

For the converse suppose $d\mid e$ and write $p=g_d(h_d)$, where $g_d$ has leading coefficient $a_n$ and degree $n/d$. Put $m_e=n/e$. The formal Laurent series
\[
 V(t)=t^{e/d}\left(\frac{g_d(t)}{a_n t^{n/d}}\right)^{1/m_e}
\]
has integer exponents, leading coefficient one, and $V^{m_e}=g_d/a_n$. Its composition with $h_d$ is well defined at infinity and has leading term $x^e$. Uniqueness of a unit formal root gives
\[
 V(h_d(x))=\left(p(x)/a_n\right)^{1/m_e}.
\]
Every negative power of $h_d$ has only strictly negative powers of $x$ in its Laurent expansion at infinity. Every positive power of $h_d$ is a polynomial with zero constant term. Therefore the strictly positive-power part of $V(h_d)$ is obtained by taking the strictly positive-power polynomial part $q$ of $V$ and substituting $h_d$. By the Laurent-series formula in Section~\ref{sec:candidate}, that positive-power part is $h_e$. Thus $h_e=q(h_d)$, as required. The argument includes $d=1$ and $e=n$.
\end{proof}

In particular one can summarize all normalized right components by a finite divisibility poset on at most $\tau(n)$ vertices. Its cover relations encode indecomposable quotients.

\begin{proposition}[Complete chains as maximal divisor chains]\label{prop:maxchains}
Let
\[
 1=d_0\mid d_1\mid\cdots\mid d_r=n
\]
be a strict chain in $S(p)$ and write $h_{d_i}=q_i\circ h_{d_{i-1}}$. Then
\begin{equation}
 p=(a_nq_r+p(0))\circ q_{r-1}\circ\cdots\circ q_1.\label{eq:posetchain}
\end{equation}
This is a complete decomposition exactly when every adjacent relation is a cover relation in the divisibility poset. All normalized complete decompositions arise this way.
\end{proposition}
\begin{proof}
Composing the quotient identities yields $h_n=q_r\circ\cdots\circ q_1$, which proves \eqref{eq:posetchain}. Suppose some quotient $q_i=A\circ B$ is decomposable. Normalize $B$; since $q_i$ is monic with zero constant term, the adjusted $A$ can also be taken monic with zero constant term. Then $B(h_{d_{i-1}})$ is a normalized right component of $p$ whose degree lies strictly between $d_{i-1}$ and $d_i$ and divides $d_i$. Thus the relation is not a cover.

Conversely, an intermediate degree $e\in S(p)$ with $d_{i-1}\mid e\mid d_i$ gives two quotient identities by Theorem~\ref{thm:poset}; their composition factors $q_i$ nontrivially. The outer affine map in \eqref{eq:posetchain} does not change decomposability. Finally the suffix compositions of any normalized complete chain are normalized right components of $p$ and provide its associated maximal divisor chain.
\end{proof}

This yields an alternative architecture: compute $S(p)$ once, recover quotient labels on its cover edges, and enumerate maximal paths. The live packages use the recursive algorithm of Theorem~\ref{thm:enumeration}; they do not expose a poset or graph API. The poset description is included as a proved structural result and a concrete route to a more efficient future enumerator.

\subsection{A field-tower interpretation}
For any nonconstant polynomial $h$ of degree $d$,
\begin{equation}
 [K(x):K(h(x))]=d.\label{eq:fielddegree}
\end{equation}
Here is an elementary proof. The element $x$ satisfies $h(T)-h(x)=0$ over $K(h(x))$, so the degree is at most $d$. If $1,x,\ldots,x^{d-1}$ were linearly dependent over $K(h)$, clearing denominators would give a nonzero relation $\sum_{i=0}^{d-1}A_i(h)x^i=0$ with $A_i\in K[t]$. The nonzero summands have distinct leading degrees modulo $d$, so their highest terms cannot cancel. This contradiction gives the opposite inequality.

A complete decomposition thus produces a tower
\[
 K(p(x))\subset K(f_2\circ\cdots\circ f_r(x))
 \subset\cdots\subset K(f_r(x))\subset K(x),
\]
whose successive degrees are the component degrees. Our algorithms do not construct splitting fields, Galois groups, or these intermediate fields explicitly.

\section{Nonuniqueness and the Ritt perspective}
\label{sec:ritt}
Fixed-degree uniqueness does not say that all decompositions have the same
ordered degrees. Ritt's characteristic-zero theory describes how complete
chains can differ. In particular, their lengths and multisets of component
degrees agree, while their order can change. The discussion of Zieve and
M\"uller provides a detailed treatment and modern perspective \cite{ZM}.
None of the proofs of the algorithms above depends on invoking Ritt's
classification as a black box.

One consequence of that cited theorem connects completeness with optimization:
every complete chain minimizes the largest nonlinear component degree.
Indeed, refining a component of degree \(ab\), \(a,b>1\), replaces it by
components of smaller degrees \(a,b\), so refinement cannot increase the
maximum.  Ritt's invariant makes the maximum identical for every complete
refinement.  This does not imply that all indecomposable degrees are prime:
the nested example below contains an indecomposable quartic.

\subsection{Powers}
For $x^{12}$, the normalized complete chains are exactly
\[
 (x^3,x^2,x^2),\qquad (x^2,x^3,x^2),\qquad (x^2,x^2,x^3).
\]
For $x^{30}$ there are six chains, corresponding to permutations of degrees
$2,3,5$. Composing two powers multiplies their exponents:
\begin{equation}
 x^r\circ x^s=x^{rs},\qquad\text{not }x^{r+s}.
 \label{eq:powerproduct}
\end{equation}
Even the correct multiplication should not be used to compress a
\emph{complete} chain, because it merges two atoms into a decomposable factor.

In general, if \(n=\prod_i p_i^{e_i}\) and \(r=\sum_i e_i\), every normalized
right component of \(x^n\) is \(x^d\) for some divisor \(d\mid n\).  Thus its
complete chains are the distinct orders of the multiset of prime factors,
numbering \(r!/\prod_i e_i!\).  For \(n=6,12,16,30,36,72\), the counts are
\(2,3,1,6,6,10\).  These are exact enumeration benchmarks independent of
another computer algebra system's decomposition routine.

\subsection{A Chebyshev collision in the package's normalization}
The identities $T_2\circ T_3=T_3\circ T_2=T_6$ yield two normalized pairs for
\[
 T_6(x)=32x^6-48x^4+18x^2-1:
\]
\begin{align*}
 T_6&=(32t^3-48t^2+18t-1)\circ x^2,\\
 T_6&=(32t^2-1)\circ\left(x^3-\frac34x\right).
\end{align*}
Their right degrees are different, so there is no conflict with
Corollary~\ref{cor:unique}. These are elementary test cases for all-chain
enumeration; an implementation that always returns one degree pattern has
not enumerated every chain.

More generally, every divisor \(d\mid n\) gives the normalized right component
\[
 H_d(x)=\frac{T_d(x)-T_d(0)}{2^{d-1}}
\]
of \(T_n\).  The identity \(T_a\circ T_b=T_{ab}\) follows by evaluating at
\(x=\cos\theta\), where both sides equal \(\cos(ab\theta)\), and then using
equality on infinitely many points.  All divisor degrees occur, so the
right-component poset and complete-chain count coincide with those of \(x^n\).

More general collisions are not merely commutations of powers. For example,
whenever the expressions have the required nontrivial degrees,
\[
 x^m\circ\bigl(x^rR(x^m)\bigr)
 =\bigl(x^rR(x)^m\bigr)\circ x^m.
\]
The package need not recognize such templates in advance: both degree tests
succeed directly.

\section{Exact positive and negative certificates}
\label{sec:certificates}
A composition identity is a useful positive check, but it does not establish
that a failed search was exhaustive or that returned factors are atoms. The
package therefore exposes the evidence behind every proposed right degree.

A fixed-degree record contains $d$, $m$, the forced $h_d$, the digit list
$r_0,\ldots,r_m$, the constant-digit outer candidate $F_d$, a Boolean outcome,
a possible obstruction, and the residual $p-F_d(h_d)$. An obstruction is a
triple $(j,i,c)$ with $i>0$ and
\[
 c=[x^i]r_j\ne0.
\]
Digit and power indices in the software record are \emph{zero-based}
mathematical indices, although Wolfram list positions are one-based. The
implementation selects the first nonzero nonconstant coefficient by increasing
digit index and then increasing power.

\begin{proposition}[A certificate checker independent of the search]
\label{prop:checker}
A fixed-degree record is sufficient to certify the outcome if a checker
verifies all of the following: $d$ is an admissible divisor; $h_d$ is monic
of degree $d$ and has constant zero; the leading congruence
\eqref{eq:topcongruence} holds; every digit has degree less than $d$; the
reconstruction \eqref{eq:digits} equals $p$; and the recorded status agrees
with whether all digits are constant. On rejection, the claimed obstruction
must be an actual nonzero nonconstant digit coefficient.
\end{proposition}
\begin{proof}
The leading congruence and normalization identify the unique possible inner
component by Lemma~\ref{lem:triangular}. Digit reconstruction and degree bounds
identify the unique polynomial-base expansion by Theorem~\ref{thm:digits}.
Its constant or nonconstant nature decides existence. Thus neither the
recurrence used to find $h_d$ nor the long-division search must be trusted as
part of this checker.
\end{proof}

The delivered verifier follows this strategy: it raises $h_d$ to $m$ by
polynomial multiplication and reconstructs the digit expression by Horner
composition. It does not call the candidate recurrence, the monic-division
routine, or the pair search. It additionally checks that the reported outer
candidate is exactly the list of digit constants and that the residual is
correct.

An exhaustive record includes a record for \emph{every} proper divisor in
increasing order. The verifier checks this entire list, including the absence
of omissions and duplicates, and checks the accepted-degree summary. If all
tests are negative, it certifies indecomposability. For prime degree the
empty divisor list already suffices. To certify a complete returned chain,
verify its composition identity and exhaustively certify every nonlinear
component as indecomposable. All this can be done through the public API.

\paragraph{Trust boundary.}
These are mathematical certificates checked by exact computer algebra, not
formally verified proof objects in a proof assistant. The search and checker
still share scalar algebraic arithmetic and some basic coefficient-vector
operations. They therefore do not provide implementation independence against
every possible shared bug. The Python implementation uses a different exact number-field
representation; the archived checks also compare the recurrence with an
independent finite-binomial construction. Current native Wolfram execution
and independent checks are reported in Section~\ref{sec:verification}.

\section{Worked examples}
\label{sec:examples}
\subsection{A complete quartic criterion}
Let
\[
 p=a_4x^4+a_3x^3+a_2x^2+a_1x+a_0,\qquad a_4\ne0.
\]
There is only one proper degree, $d=2$. Define
\[
 b=\frac{a_3}{2a_4},\qquad h=x^2+bx,\qquad
 c=a_2-a_4b^2.
\]
The recovered outer polynomial and residual are
\[
 G(t)=a_4t^2+ct+a_0,\qquad R=(a_1-bc)x.
\]
Therefore
\begin{equation}
 \boxed{p\text{ decomposes}\ \Longleftrightarrow\
 8a_4^2a_1-4a_4a_3a_2+a_3^3=0.}\label{eq:quartic-criterion}
\end{equation}
This is valid for arbitrary algebraic coefficients, including nonreal coefficients and nonradical \code{Root} values. It follows by direct expansion, so no polynomial factorization or algebraic equation solving is hidden in the criterion.

\subsection{Both possible sextic degree patterns}
For $p=\sum_{j=0}^6a_jx^j$, $a_6\ne0$, the only possible proper right degrees are $2$ and $3$. The following outer polynomials and residuals use descending-pivot recovery; on rejection they need not equal the live digit-constant certificate fields.

For $d=2$, put
\[
 b=\frac{a_5}{3a_6},\qquad h=x^2+bx,
\]
\[
 C=a_4-3a_6b^2,\qquad D=a_2-Cb^2.
\]
Then
\begin{align*}
 G(t)&=a_6t^3+Ct^2+Dt+a_0,\\
 R(x)&=(a_3-a_6b^3-2Cb)x^3+(a_1-Db)x.
\end{align*}
Thus both displayed residual coefficients must vanish.

For $d=3$, define
\[
 u=\frac{a_5}{2a_6},\qquad
 v=\frac{a_4}{2a_6}-\frac{a_5^2}{8a_6^2},\qquad h=x^3+ux^2+vx,
\]
and $B=a_3-2a_6uv$. Then
\begin{align*}
 G(t)&=a_6t^2+Bt+a_0,\\
 R(x)&=(a_2-a_6v^2-Bu)x^2+(a_1-Bv)x.
\end{align*}
Again the two coefficients give a necessary and sufficient test. Testing only the quadratic inner candidate would miss valid degree pattern $(2,3)$; testing only the cubic inner candidate would miss pattern $(3,2)$.


\subsection{The original radical polynomial}
Put $s=\sqrt2$ and take the polynomial in the question:
\begin{align}
 p(x)={}&3+3s+(14+4s)x+(12+26s)x^2\notag\\
        &+(56+8s)x^3+(8+48s)x^4+48x^5+16s x^6.
 \label{eq:original}
\end{align}
For $d=2$, $m=3$ and $s_1=48/(16s)=3/s$. Hence $u_1=1/s$ and
\begin{equation}
 h_2=x^2+\frac{x}{s}.
 \label{eq:originalh}
\end{equation}
All base-$h_2$ digits are constant, giving
\begin{equation}
 F_2(t)=3+3s+(8+14s)t+(8+24s)t^2+16s t^3.
 \label{eq:originalF}
\end{equation}
Direct substitution verifies $p=F_2(h_2)$. This differs from the question's
pair only by the affine normalization $g=s h_2$:
\[
 g=x+s x^2,\qquad
 f(t)=3+3s+(14+4s)t+(4+12s)t^2+8t^3,
\]
and indeed $F_2(t)=f(st)$.

The alternative right degree $d=3$ produces
\[
 h_3=x^3+\frac{3s}{4}x^2+\left(\frac{15}{16}+\frac{s}{8}\right)x.
\]
Its digit expansion is $p=r_0+r_1h_3+r_2h_3^2$ with
\begin{align*}
 r_0={}&3+3s+\left(\frac{51}{16}+\frac{3s}{4}\right)x
                +\left(\frac32+\frac{51s}{16}\right)x^2,\\
 r_1={}&11+2s,\qquad r_2=16s.
\end{align*}
The nonconstant $r_0$ proves that this degree is impossible. Thus the original
polynomial has exactly one normalized complete chain, of degrees $(3,2)$.

\begin{lstlisting}
p = 3 + 3 Sqrt[2] + (14 + 4 Sqrt[2]) x +
    (12 + 26 Sqrt[2]) x^2 + (56 + 8 Sqrt[2]) x^3 +
    (8 + 48 Sqrt[2]) x^4 + 48 x^5 + 16 Sqrt[2] x^6;
chain = AlgebraicDecompose[p, x];
VerifyAlgebraicDecomposition[p, chain, x]
AlgebraicRightDecompose[p, x, 3]
(* Mathematical expectations: True; Missing["NotDecomposable",3]. *)
\end{lstlisting}
The displayed formulas are exact mathematical results independently checked
in the bundled verifier; they are not represented as a native Wolfram session
transcript.

\subsection{The nested degree-24 example}
Let $p_2(x)=p(x^4-x+1)$ and set $v=x^4-x$. Since
\[
 h_2(v+1)=v^2+\left(2+\frac1s\right)v+1+\frac1s,
\]
a normalized chain is
\begin{align*}
 H_3(x)&=x^4-x,\\
 H_2(x)&=x^2+\left(2+\frac1s\right)x,\\
 H_1(x)&=141+105s+(168+142s)x\\
       &\hspace{12mm}+(56+72s)x^2+16s x^3.
\end{align*}
Thus $p_2=H_1\circ H_2\circ H_3$. The quartic $H_3$ is indecomposable: its
only proposed right degree is two, the leading equations force $x^2$, and
the term $-x$ obstructs membership. The quadratic and cubic are
indecomposable by prime degree. The returned degree sequence is $(3,2,4)$.
In particular, a component of composite degree can be an atom.

\subsection{An algebraic \code{Root} coefficient}
Let $\alpha$ be the particular real algebraic number selected in Wolfram
notation by
\begin{lstlisting}
a = Root[#^5 - # - 1 &, 1];
\end{lstlisting}
and take
\[
 q(x)=(x^3+\alpha x)^2+(1+\alpha)(x^3+\alpha x)+2.
\]
For right degree three the algorithm recovers exactly
\[
 (t^2+(1+\alpha)t+2,\quad x^3+\alpha x).
\]
No radical formula for $\alpha$ is requested or constructed. The coefficient
relation $\alpha^5-\alpha-1=0$ is available to exact algebraic reduction.
The same construction works for a nonreal root of this polynomial and for
polynomials with coefficients from several algebraic extensions at once.

A useful validation example adds $(\alpha^5-\alpha-1)x^{40}$ to $x^4$.
The true polynomial is still $x^4$, so its degree must be recomputed after
coefficient reduction. The package does this before selecting divisors; it
does not trust the apparent degree of unreduced input syntax.

\subsection{Complex coefficients and affine normalization}
Over $\Q(i,\sqrt2)$,
\[
 (1+i)(x^3+ix)^2+\sqrt2
 =\bigl((1+i)t^2+\sqrt2\bigr)\circ(x^3+ix).
\]
For a less normalized input, choose
\[
 f(t)=(1+\sqrt2)t^2+(2-\sqrt3)t+7,\qquad
 g(x)=(2+\sqrt3)x^3+\sqrt2 x+\sqrt3.
\]
The unique normalized degree-three pair is
\[
 \left(f((2+\sqrt3)t+\sqrt3),\quad
           x^3+\frac{\sqrt2}{2+\sqrt3}x\right).
\]
These examples emphasize that the outer leading coefficient, the original
inner leading coefficient, and the original inner constant term are all
handled without ad hoc assumptions.

\subsection{Composite-degree atoms and exact perturbations}

For every \(n\ge2\) and nonzero \(c\), the polynomial \(x^n+cx\) is
indecomposable over every characteristic-zero extension.  Prime degree is
immediate.  Otherwise every proper right degree satisfies \(d\le n/2\);
the leading equations force \(h_d=x^d\), and the digit \(r_0=cx\) is
nonconstant.  The fixed-degree theorem excludes every possible decomposition.
Thus \(x^n+10^{-100}x\) is exactly indecomposable although arbitrarily close
to \(x^n\).  Approximate agreement is not an exact zero certificate.

\section{Comparison and reconciliation of the three reports}\label{sec:reports}

The three original distributions are preserved byte-for-byte under
\path{polynomial-decompose/reports/}; the live code and this article are
separate derived work.  All three develop the same normalized-candidate
theorem, formal-root recurrence, descent, and finite decomposition algorithms.
Their differences matter for implementation and certification.

\begin{table}[htbp]
\centering\small
\begin{tabularx}{\textwidth}{@{}lX@{}}
\toprule
Report & Distinctive contributions\\
\midrule
\cite{r1} & \code{AlgebraicDecomposition-1.0.0.zip}: monic-division digits,
explicit normalization and descent, recursive refinement of both factors,
finite enumeration limits, the Ritt minimum-maximum-degree consequence,
and the careful order-dependent analysis of derivative recursion.
Its independent checks diagnose a missed-decomposition case in the
inspected SymPy 1.14.0 implementation.\\
\cite{r2} & \code{algebraic-polynomial-decomposition.zip}: descending-pivot
outer recovery and residual certificates, explicit quartic and sextic
criteria, a self-contained proof of the right-component divisibility
poset, and the smooth fixed-degree decomposition locus.  Its diagnostic
and verification interfaces differ from the digit-based design.\\
\cite{r3} & \code{algebraic\_polynomial\_decomposition.zip}: the digit-certificate
API adopted by the live implementation, an independent certificate checker
using top congruence and digit reconstruction, deterministic
smallest-right-degree recursion, and extensive exact-domain comparisons.\\
\bottomrule
\end{tabularx}
\caption{The reports agree on the core mathematics but differ in recovery,
certificates, enumeration, and software interfaces.}\label{tab:reports}
\end{table}

Two mathematical reconciliations are explicit in this article.  On rejection,
the digit and pivot outer candidates are not interchangeable; their different
residuals are illustrated in Section~\ref{sec:digits}.  Also, outer-only
derivative recursion is not inherently incorrect: report~\cite{r1} supplies
the precise candidate-order condition under which it is justified
(Section~\ref{sec:audit}).

All three reports stated that their Wolfram code was unexecuted in the
preparation environment.  Current native reruns show why that limitation
matters.  Report~1 passes \(55/107\) tests, report~2 passes \(39/53\), and
report~3 passes \(60/60\) under Wolfram 15.0.1.  The first two return
\(\{x^4\}\) as a purported complete decomposition of \(x^4\), while their
pair routines find \(\{x^2,x^2\}\); report~2's own completeness verifier also
incorrectly accepts the singleton.  The diagnosed pair-search control flow
uses \code{Return} inside a \code{Do} loop as though it exited the enclosing
function.  The live code adopts report~3's explicit loop-exit structure.
These failures concern native implementation, not counterexamples to the
mathematical theorem.  Independent archived Python checks continue to pass;
their narrower evidentiary scope is recorded in Section~\ref{sec:verification}.

\section{Wolfram Language and Python implementation}
\label{sec:implementation}

\subsection{Loading and basic operation}
Load the live Wolfram file from the project directory, with an unassigned
polynomial variable:
\begin{lstlisting}
Get["AlgebraicDecomposition.wl"];
Clear[x];
AlgebraicDecompose[(x^4 + x)^2, x]
(* {x^2, x^4 + x}, up to exactly equal formatting *)
\end{lstlisting}
All factor lists are outermost first.  Every component except the first is
monic with zero constant term.  Neither coefficient formatting nor a choice
between radicals and \code{Root} objects is part of the equality contract.
The package does not require Python or an internet connection for evaluation.

\subsection{Coefficient arithmetic and representation}
The implementation uses ascending dense coefficient vectors. For example,
$3+2x+x^3$ is represented internally by $\{3,2,0,1\}$. The zero polynomial
has the unique vector $\{0\}$. Trailing zero entries are removed only after
exact scalar normalization.

The scalar normalization boundary is the private function \code{red}. Exact
integers and rationals pass directly. Other coefficients are processed with
\code{RootReduce}, then checked for exact numeric membership in
\code{Algebraics}. Wolfram documents \code{RootReduce} as reducing combinations
of exact algebraic quantities to a single algebraic-number representation,
and \code{Algebraics} provides the corresponding exact domain membership
operation \cite{RootReduce,Algebraics}. The algorithm applies these operations
to coefficients, not to polynomials containing the free variable.

An approximate real anywhere in the evaluated coefficient expression is
rejected. So are nonnumeric symbolic parameters and quantities not certified
as algebraic. A failed exact reduction returns a \code{Failure}; it is never
treated as a nonzero coefficient or as evidence of indecomposability. After
normalization, equality to the exact integer zero is used for zero tests.
There is no tolerance, numerical sampling, or \code{PossibleZeroQ} in the
algorithm.

This contract covers algebraic \code{Root} objects and exact arithmetic with
them, not arbitrary \code{Root} expressions encoding transcendental equations.
It also does not promise to recognize every sophisticated expression that
happens mathematically to evaluate to an algebraic number. An explicit
algebraic representation can be supplied when the initial expression lies
outside the recognized domain.

\subsection{Public functions}
All lists of factors are outermost first. Mathematical equality, rather than
the printed choice between radicals and \code{Root} objects, is the output
contract.

\paragraph{\code{AlgebraicDecompose[p,x]}.}
Returns one complete normalized chain. The smallest successful right degree
is chosen at each step, giving deterministic degree choices. An input of
degree at most one returns its singleton list.

\paragraph{\code{AlgebraicDecompositions[p,x]}.}
Returns all normalized complete chains, modulo affine changes at internal
interfaces. It implements \eqref{eq:enumeration}, not an enumeration of all
raw affine variants or all incomplete chains. The output can be large.
The option \code{"MaxDecompositions" -> Infinity} sets the default unlimited
enumeration. A positive integer limit is exceeded only when a further chain
is actually found. In that event the result is
\code{Failure["EnumerationLimit",data]}, with \code{"Complete" -> False},
the requested \code{"Limit"}, and \code{"PartialDecompositions"}. Exhausting
the search with exactly the limit returns the complete list normally; a
truncated list is never silently presented as exhaustive.

\paragraph{\code{AlgebraicDecompositionPairs[p,x]}.}
Returns all normalized nontrivial pairs, ordered by increasing right degree.
An empty list means no nontrivial pair exists for a valid input. This interface
is often a better diagnostic than full chain enumeration.

\paragraph{\code{AlgebraicRightDecompose[p,x,d]}.}
Returns the unique normalized pair for a specified proper divisor $d$.
A valid degree with no decomposition returns
\code{Missing["NotDecomposable",d]}. An invalid degree, such as a nondivisor,
one, or the whole input degree, returns \code{Failure}. This distinction
separates a proved negative answer from an invalid request.

\paragraph{\code{AlgebraicDecompositionData[p,x,d]}.}
Returns the fixed-degree certificate record. Its keys are \code{"Type"},
\code{"RightDegree"}, \code{"OuterDegree"}, \code{"Inner"},
\code{"OuterCandidate"}, \code{"Digits"}, \code{"Decomposable"},
\code{"Obstruction"}, and \code{"Residual"}. The outer candidate is not a
valid factor when the status is false. The obstruction is \code{None} on
success, otherwise an association with \code{"DigitIndex"}, \code{"Power"},
and \code{"Coefficient"}.

\paragraph{\code{AlgebraicDecompositionData[p,x]}.}
\begingroup\raggedright
Returns an exhaustive report with \code{"InputDegree"},
\code{"TestedRightDegrees"}, \code{"AcceptedRightDegrees"},
\code{"Indecomposable"}, and \code{"Tests"}, in addition to its type marker.
For degree below two, the indecomposability field is
\code{Missing["NotApplicable","DegreeBelowTwo"]}.\par\endgroup

\paragraph{\code{VerifyAlgebraicDecompositionData[p,data,x]}.}
Checks a fixed-degree or exhaustive report using
Proposition~\ref{prop:checker}. Malformed or inconsistent data returns
\code{False}; invalid polynomial input returns \code{Failure}. The checker
is intended for mathematical data, not as a sandbox for evaluating untrusted
Wolfram expressions.

\paragraph{Composition and chain verification.}
\code{ComposeDecomposition[parts,x]} composes a list exactly; the empty list
represents the identity $x$.
The default verification call is
\begin{lstlisting}
VerifyAlgebraicDecomposition[p, parts, x]
\end{lstlisting}
It checks the composition identity. For example, $\{x^4\}$ passes that check for $p=x^4$ but
is not a complete nonlinear chain. The options
\code{"RequireComplete" -> True} and \code{"RequireNormalized" -> True}
add exhaustive atomicity and internal normalization checks, respectively;
both default to \code{False}. The singleton conventions for inputs of degree
at most one are recognized by the completeness option.

\subsection{Checking a complete chain with public operations}
The following pattern distinguishes the identity check from the atomicity
checks:
\begin{lstlisting}
chain = AlgebraicDecompose[p, x];
identityOK = VerifyAlgebraicDecomposition[p, chain, x];
reports = AlgebraicDecompositionData[#, x] & /@ chain;
certificatesOK = And @@ MapThread[
  VerifyAlgebraicDecompositionData[#1, #2, x] &, {chain, reports}];
atomsOK = AllTrue[reports, TrueQ[#["Indecomposable"]] &];
(* For degree(p)>=2, require identityOK && certificatesOK && atomsOK. *)
\end{lstlisting}
In production code, check for \code{Failure} before inspecting returned lists
or associations. Inputs should use an unassigned symbol as the polynomial
variable; ordinary Wolfram evaluation happens before the function patterns
are matched.

\subsection{The Python interface}
The independent live module is \code{python/algebraic\_decompose.py}.
Its public functions mirror the mathematical operations:
\begin{lstlisting}
decompose(p, x)
decompositions(p, x, *, max_chains=None)
decomposition_pairs(p, x)
right_decompose(p, x, d)
decomposition_data(p, x, d=None)
verify_decomposition_data(p, data, x)
compose(parts, x)
verify_decomposition(p, parts, x, *,
                     require_complete=False, require_normalized=False)
\end{lstlisting}
The optional limits and verifier flags are keyword-only parameters.
A valid negative \code{right\_decompose} returns \code{None}; invalid requests
raise \code{ValueError}. Certificate keys use the corresponding snake-case
names. The default \code{max\_chains=None} is unlimited. Finding a chain
beyond a positive limit raises \code{EnumerationLimitError}, whose
\code{limit}, \code{partial\_chains}, and \code{complete=False} attributes
make the incomplete result explicit. Exactly attaining a limit is successful
when traversal is then exhausted.

An exact SymPy expression or a univariate \code{Poly} over \code{QQ} or an
\code{AlgebraicField} is accepted. The implementation chooses one exact
coefficient field per call, caches scalar conversions, and performs the
polynomial algorithms in that field. Optional FLINT support accelerates
rational multiplication and division; it does not change the criterion.
A coefficient represented by \code{CRootOf} is treated as a scalar even
when its defining polynomial uses the same printed variable as the input.
Approximate coefficients, unbound parameters, and unsupported exact
transcendental coefficients are rejected. Both implementations verify the
selected exact algebraic numbers, without numeric tolerance.

\subsection{Algorithmic organization}
The fixed-degree engine consists of three small components: the formal-root
recurrence, monic long division, and detection of nonconstant digits. Outer
coefficients are taken directly from digit constants. Composition and
certificate checking use explicit polynomial convolution and Horner schemes.
Loop exits are explicit: the nested obstruction scan uses a locally tagged
\code{Catch}/\code{Throw}, and pair search records its result before
\code{Break}. Wolfram's \code{Return} inside a \code{Do} loop exits that loop,
not necessarily the surrounding function, so it must not be used as an
implicit multilevel exit \cite{Return}.
There is no call to \code{Factor}, \code{Decompose}, \code{Solve}, or a radical
denesting routine in the search. The complete live sources are the loadable
\code{AlgebraicDecomposition.wl} and
\code{python/algebraic\_decompose.py}; they are not duplicated as a source appendix,
which could diverge from the tested files.

There are deliberately no speculative options for extensions, numerical
tolerances, or heuristic fallback solvers. Unsupported input or a failed
algebraic normalization does not silently change the mathematical problem.
The implementation also does not automatically convert every output to
radicals: such a conversion may be impossible or needlessly expensive, and
would not improve correctness.

\section{Complexity, coefficient growth, and practical limits}
\label{sec:complexity}
\subsection{Field-operation bounds}
For a fixed proposed right degree $d$, the recurrence performs
$1+2+\cdots+(d-1)=O(d^2)$ scalar operations. A dense division of a polynomial
of degree $N$ by a monic degree-$d$ polynomial costs $O((N-d+1)d)$ field
operations. Applying this repeatedly at degrees $n,n-d,n-2d,\ldots$ costs
\[
 O\!\left(d\sum_{j=0}^{m-1}(n-jd)\right)=O(n^2),
 \qquad m=n/d.
\]
Thus the fixed-degree test is $O(n^2)$ in the dense field-operation model,
with $O(n)$ field elements of working storage for its coefficient vectors and
digits. This agrees with the classical quadratic fixed-degree approach
\cite{KL}; it is not a bit-complexity bound.

Testing every proper divisor gives $O(\tau(n)n^2)$ field operations. Along
one complete chain, subsequent outer degrees are divisors of $n$ and at most
half the previous degree. Since $\tau(n_j)\le\tau(n)$,
\[
 \sum_j O(\tau(n_j)n_j^2)
 \le O\!\left(\tau(n)n^2\sum_{j\ge0}4^{-j}\right)
 =O(\tau(n)n^2).
\]
Enumeration of all chains adds the cost of its search graph and its output;
no polynomial bound in input size alone is claimed for that interface.
Storing all degree certificates simultaneously also costs more than storing
one fixed-degree test.

\subsection{Why algebraic arithmetic is not constant-cost}
Let $D=[K:\Q]$. In a common primitive-element representation $K=\Q(\theta)$,
an element is a rational vector of length $D$, interpreted modulo the minimal
polynomial of $\theta$. Multiplication, inversion, and rational height growth
all have costs depending on $D$ and coefficient sizes. A concise expression
with many independent radicals can generate a field whose degree is the
product of their individual degrees. Therefore a polynomial field-operation
count does not imply polynomial running time in the length of an arbitrary
radical expression.

The delivered Wolfram backend normalizes the input coefficients once and
uses a rational fast path, but nonrational arithmetic still calls
\code{RootReduce}. It may repeat expensive algebraic work, especially when
many unrelated generators occur. A future Wolfram backend could convert all
coefficients to one common field and operate on rational coordinate vectors;
\code{ToNumberField} is relevant to such a design \cite{ToNumberField}.
The live Python implementation already uses one explicit exact coefficient
field, with cached scalar conversions and optional FLINT acceleration for
rational polynomial arithmetic. Python is not a runtime dependency of the
Wolfram package.

\subsection{Storage, sparsity, and safe early rejection}
Both implementations use dense coefficient vectors. A sparse expression such
as $x^{1000000}+x$ still creates a long vector. Within that representation,
the live engines avoid unnecessary work: if the leading $d-1$ normalized
coefficients vanish, the candidate is immediately $x^d$; monic division skips
zero coefficients; and ordinary searches stop at the first nonconstant digit.
The certificate interface continues division to retain the full
reconstruction. These optimizations preserve Theorem~\ref{thm:digits}.

The descending-pivot method likewise costs $O(n^2)$ field operations with
ordinary dense arithmetic: successive powers of $h$ cost
$O(d^2m^2)=O(n^2)$, and the eliminations cost $O(nm)$.
Storing every power simultaneously requires $O(n^2/d)$ field elements;
the digit method uses $O(n)$ working elements for a fixed degree.
Retaining all degree certificates and all output chains has additional
storage costs. Large exact coefficient heights and the common field degree
can dominate either algorithm. A native measurement of one sparse rejection
case appears in Section~\ref{sec:verification}; it is not a general speed guarantee.

\section{Audit of the derivative-factorization approach}
\label{sec:audit}
The identity
\[
 p'(x)=f'(g(x))g'(x)
\]
is mathematically useful. In characteristic zero, after suitable scalar
normalization, the derivative of a right component divides $p'$. An exhaustive
enumeration of derivative divisors, followed by integration, normalization,
and exact membership testing, can therefore produce a correct algorithm.
The factorization-free method replaces this potentially large candidate set
by one forced candidate per right degree.

Several details in the posted code require correction independently of its
choice of method.

\paragraph{One distinct derivative factor is not an obstruction.}
The early return when the nonconstant factor list has length below two is
incorrect. For $p=x^4$, $p'=4x^3$ has a single distinct nonconstant factor
$x$ with multiplicity three, yet $p=x^2\circ x^2$. Multiplicity cannot be
ignored when deciding whether candidate derivatives exist.

\paragraph{Repeated-factor subsets create duplicates.}
Expanding multiplicities into repeated copies and taking all subsets can
visit $2^{\sum e_i}$ subsets, with many yielding the same product. Divisors
of $\prod q_i^{e_i}$ are more naturally indexed by multiplicity tuples
$0\le j_i\le e_i$, of which there are $\prod(e_i+1)$. Even that is an
avoidable combinatorial search for the present characteristic-zero task.

\paragraph{The candidate degree must divide the input degree.}
If a proposed derivative has degree $d-1$, then $d$ must be a proper divisor
of $n$. The posted degree bound on derivative products does not enforce this.
Using $n/d$ subsequently as the number of outer coefficients can therefore
produce a noninteger degree and an invalid coefficient system. Divisibility
must be checked before constructing the outer polynomial.

\paragraph{Outer-only recursion has an order-dependent justification.}
A successfully found inner factor is not automatically an atom.  However,
report~\cite{r1} correctly observes that an exhaustive derivative search
in increasing number of irreducible-factor copies can justify the posted
outer-only recursion.  If the selected \(h\) had \(h=A\circ B\), then
\(B'\) would be a proper divisor of \(h'=(A'\circ B)B'\), with fewer
factor copies.  Its integrated candidate would be a right component of
the input and would succeed earlier under complete exact membership testing.
Under those conditions, the first successful inner component is indecomposable.

This is a proof obligation on the actual ordering and completeness of the
tests.  Search caps, a changed order, or an incomplete solver invalidate it.
The live single-chain algorithm instead uses increasing right degree and
Lemma~\ref{lem:minimal}, which gives the atom guarantee explicitly.

\paragraph{The final power rewrite changes the composition.}
The replacement of adjacent powers by $x^{n+m}$ is algebraically wrong;
\eqref{eq:powerproduct} requires $x^{nm}$. For example, adjacent cubic powers
compose to $x^9$, not $x^6$. Correcting addition to multiplication fixes the
identity but still destroys completeness by merging nonlinear factors.

\paragraph{Exact checking must remain exact.}
A heuristic possible-zero predicate is not the same interface contract as
coefficientwise exact algebraic reduction. Likewise, the existence of one
solver result is not an exhaustive negative certificate when no result is
returned. For explicit algebraic coefficients, reduction to canonical exact
numbers and the finite obstruction test make the decision boundary explicit.

The historical comment in the source thread describes an old implementation
context for \code{Decompose}; it does not establish that the derivative
approach is the only known algorithm. The factorization-free leading-term
approach substantially predates that discussion \cite{KL,SE}.

\section{Validation record and reproducibility}\label{sec:verification}

\subsection{Current native execution}
The original native suites were run from temporary copies, preserving all
report files.  The runtime was Wolfram 15.0.1 on 64-bit Windows.
\begin{center}
\begin{tabular}{@{}lrrl@{}}
\toprule
Package & Passed & Total & Result\\
\midrule
Original report 1 & 55 & 107 & 51 wrong-result failures, 1 message failure\\
Original report 2 & 39 & 53 & 14 wrong-result failures\\
Original report 3 & 60 & 60 & Successful\\
Live package & 80 & 80 & Successful\\
\bottomrule
\end{tabular}
\end{center}
The live suite retains the original report~3 fixtures and adds checks for
output limits, certificate corruption, completeness and normalization
options, and input contracts.  Passing the independent mathematical checks
of the first two reports did not prevent their native control-flow failures.
A lexical bracket/comment/string scan is likewise not a native execution
receipt.

\subsection{Independent archived checks}
All three archived Python suites were rerun successfully on copies, using
Python 3.14.4 and SymPy 1.14.0:
\begin{itemize}
\item Report~1 completed 574 checks with no failures, including direct
division and finite-binomial candidate comparisons.
\item Report~2 completed 144 generated composites, 72 independent triangular
candidate comparisons, 348 residual checks, 27 indecomposable-family tests,
28 external rational comparisons, and 16 examples or boundary cases.
\item Report~3 covered 974 polynomials, 2,001 fixed-degree tests, 1,042 complete
chains, 95 known normalized pairs, and 848 rational comparisons with SymPy.
It records 14 differing degree patterns rather than claiming universal agreement.
\end{itemize}
These categories overlap and are not summed into a count of independent
experiments.  The archived scripts still print their historical
native-unavailable statements; the current native results above are separate
evidence.  The live cross-language harness passes all 15 common cases: it checks exact
substitution, one-chain, all-chain and pair parity, and independently checks
Python-generated exhaustive certificates in Wolfram. The cases include the
motivating sextic, a nested degree-24 input, the sparse-inner and negative-digit
regressions, monomial and Chebyshev collisions, generic and prime-degree inputs,
complex and mixed-quadratic coefficients, a quintic root whose bound variable
matches the polynomial variable, exact leading-coefficient cancellation, and
zero, constant, and linear inputs. The live Python suite passes all 16 test methods, including its generated and boundary cases. All current command lines are recorded in the project README.

A candidate formula independent of the differential recurrence is the finite
binomial expansion
\[
 (1+w)^{1/m}\equiv
 \sum_{j=0}^{d-1}\binom{1/m}{j}w^j\pmod{z^d},
 \qquad w\in zK[z].
\]
Constructed composites provide known normalized pairs without asking another
decomposition routine to invent the answer.  Monomial and Chebyshev chain
counts test enumeration independently.  Deliberate changes to candidate
coefficients, digit lists, accepted-degree summaries, and obstruction fields
test that certificate checking rejects malformed evidence.

\subsection{The external oracle is a comparison, not a completeness proof}
Reports~\cite{r1,r3} retain missed-decomposition discrepancies from their
inspected SymPy 1.14.0 installation.  A small example is
\[
 p=(x^3+x^2+x)^2.
\]
The external routine with explicit rational-field coefficients returned a
singleton, although the construction supplies a cubic inner component.
Report~1 diagnoses a missing term in the inspected candidate recurrence;
its independent finite-series calculation gives the correct candidate.
This is a historical version-specific observation, not a claim about every
current SymPy release.

Exact recomposition of a singleton always succeeds.  Consequently an external
routine's identity check cannot establish that its factors are indecomposable,
and agreement with an external singleton is not an indecomposability oracle.
The independent degree tests, reconstruction certificates, and known
constructions are the completeness checks used here.

\subsection{A measured sparse-rejection improvement}
In fresh Wolfram 15.0.1 kernels, excluding startup,
\[
 \code{AlgebraicDecompositionPairs}[x^{96}+x+1,x]
\]
returns the empty list in both report~3 and the live implementation.
A one-second \code{RepeatedTiming} measurement gave \(0.5750906\) seconds
for report~3 and \(0.02458484\) seconds for the live code, about a
\(23.4\)-fold improvement.  This specific sparse negative example benefits
from the monomial-candidate shortcut, skipping zero divisor coefficients,
and stopping at the first nonconstant digit in ordinary searches.
It is not a general speed ratio or a cross-language benchmark.

\subsection{Reproduction and trust boundaries}
From \code{polynomial-decompose/}, run the native suite with
\begin{lstlisting}
wolfram -script RunTests.wl
\end{lstlisting}
The live Python suite and the independent cross-language harness are
documented alongside \code{python/algebraic\_decompose.py}.  The original
report runners remain under their respective \code{reports/report-0i/}
directories.  Their output locations belong to the archived layouts, so
reproduction should use copies when preserving the report bytes.

Native tests establish that the specified cases execute correctly in the
reported runtime.  Independent certificates establish the represented
mathematical facts using their exact arithmetic backend.  Neither is
proof-assistant verification of every implementation path, and article PDF
rendering is a separate layout check rather than a mathematical test.

\section{Conclusions and further implementation work}
\label{sec:conclusion}
For explicit exact algebraic coefficients, functional decomposition has a
clean finite solution: normalize the inner component, force its coefficients
from the top of the input, and check the lower coefficients through polynomial
base digits. Each divisor contributes exactly one candidate. Successful
pairs descend to the original coefficient field, and rejected candidates
carry exact algebraic obstructions. Complete chains and all normalized
alternatives follow by elementary recursion.

The live implementations realize this theory without derivative factorization
or coefficient solving.  Their ordinary searches can stop on the first
obstruction, while certificate requests retain complete reconstruction data.
Caching, sparse arithmetic, and exact fixed coefficient fields improve
practical performance without weakening the decision criterion.
The right-component poset gives a proved further route to shared algebraic
work and path-based enumeration.  These are improvements to engineering
and scale, not missing cases in the characteristic-zero mathematical criterion.

\begingroup
\raggedright
\begin{thebibliography}{99}
\bibitem{r1}
\emph{Functional Decomposition of Polynomials with Algebraic Coefficients}.
Original report distributed as \path{AlgebraicDecomposition-1.0.0.zip};
preserved under \path{polynomial-decompose/reports/report-01/}.
Article: \path{article/article.tex}. Archived source, tests, and oracle notes
are retained as provenance, including their historical validation limitations.

\bibitem{r2}
\emph{Exact Functional Decomposition of Polynomials with Algebraic Coefficients}.
Original report distributed as \path{algebraic-polynomial-decomposition.zip};
preserved under \path{polynomial-decompose/reports/report-02/}.
Article: \path{article/algebraic-decomposition.tex}.
The live synthesis uses its triangular recovery, explicit criteria,
and right-component-poset arguments.

\bibitem{r3}
\emph{Exact Functional Decomposition over Algebraic Number Fields}.
Original report distributed as \path{algebraic_polynomial_decomposition.zip};
preserved under \path{polynomial-decompose/reports/report-03/}.
Article: \path{article/article.tex}.
The live implementations derive from its digit-certificate interface and
correct explicit loop-exit design.

\bibitem{SE}
V. Reshetnikov and contributors,
\emph{Is it possible to make Decompose work with coefficients containing
radicals?} Mathematica Stack Exchange, question 206618 and answer 206622,
2019. Source question, contributed code, and historical discussion.
\url{https://mathematica.stackexchange.com/questions/206618/}

\bibitem{KL}
D. Kozen and S. Landau,
\emph{Polynomial decomposition algorithms},
Journal of Symbolic Computation \textbf{7} (1989), 445--456.
\url{https://doi.org/10.1016/S0747-7171(89)80027-6}

\bibitem{Gathen}
J. von zur Gathen,
\emph{Functional decomposition of polynomials: the tame case},
Journal of Symbolic Computation \textbf{9} (1990), 281--299.
\url{https://vonzurgathen.online/old/files/gat90c.pdf}

\bibitem{WZ}
B. K. Wyman and M. E. Zieve,
\emph{Two questions on polynomial decomposition},
Quarterly Journal of Mathematics \textbf{63} (2012), 507--511;
arXiv:1301.5211.
In particular, the coefficient-descent argument in Proposition 2.1.
\url{https://arxiv.org/abs/1301.5211}

\bibitem{ZM}
M. E. Zieve and P. M\"uller,
\emph{On Ritt's polynomial decomposition theorems}, arXiv:0807.3578, 2008.
\url{https://arxiv.org/abs/0807.3578}

\bibitem{Decompose}
Wolfram Research, \emph{Decompose}, Wolfram Language documentation.
\url{https://reference.wolfram.com/language/ref/Decompose.html}

\bibitem{RootReduce}
Wolfram Research, \emph{RootReduce}, Wolfram Language documentation.
\url{https://reference.wolfram.com/language/ref/RootReduce.html}

\bibitem{Algebraics}
Wolfram Research, \emph{Algebraics}, Wolfram Language documentation.
\url{https://reference.wolfram.com/language/ref/Algebraics.html}

\bibitem{ToNumberField}
Wolfram Research, \emph{ToNumberField}, Wolfram Language documentation.
\url{https://reference.wolfram.com/language/ref/ToNumberField.html}

\bibitem{TestReport}
Wolfram Research, \emph{TestReportObject}, Wolfram Language documentation.
\url{https://reference.wolfram.com/language/ref/TestReportObject.html}

\bibitem{Return}
Wolfram Research, \emph{Return}, Wolfram Language documentation,
especially the example under ``Possible Issues.''
\url{https://reference.wolfram.com/language/ref/Return.html}

\bibitem{SymPy}
SymPy development team, \emph{Polynomial manipulation: reference},
SymPy documentation; local verification used version 1.14.0.
\url{https://docs.sympy.org/latest/modules/polys/reference.html}

\end{thebibliography}
\endgroup


\end{document}
